\section{Алгоритм Фано. Оптимальное кодирование. Алгоритм Хаффмана. Дерево Хаффмана.}

\subsection*{Оптимальное кодирование}

\paragraph{Постановка задачи}
Пусть задан алфавит $S = \{s_1, \dots, s_n\}$ и вероятности появления каждого символа $p_1, \dots, p_n$. Задача \textbf{оптимального кодирования} — построить такую схему префиксного кодирования, при которой \textbf{средняя длина кодового слова} будет минимальной:
$$L_{cp} = \sum_{i=1}^n p_i \cdot l_i \to \min,$$
где $l_i$ — длина кодового слова для символа $s_i$. Основная идея: часто встречающимся символам назначаются короткие коды, а редким — длинные.

\subsection*{Алгоритм Фано (Шеннона-Фано)}

Это «нисходящий» (top-down) метод построения кода.

\paragraph{Алгоритм:}
1. Отсортировать символы алфавита по убыванию вероятностей.
2. Разделить алфавит на две части так, чтобы суммарные вероятности обеих частей были максимально близки друг к другу.
3. Первой части присвоить в качестве первого разряда кода «0», второй — «1».
4. Рекурсивно повторить процедуру для каждой части, пока в группах не останется по одному символу.

\paragraph{Особенности:}
Метод Фано гарантирует получение префиксного кода, но не всегда дает математически оптимальный результат (иногда суммы вероятностей нельзя разделить достаточно ровно).

\subsection*{Алгоритм Хаффмана}

Это «восходящий» (bottom-up) метод, который всегда строит \textbf{математически оптимальный} префиксный код.

\paragraph{Алгоритм:}
1. Выписать все символы как свободные узлы будущего дерева, указав их вероятности (веса).
2. Выбрать два узла с \textbf{наименьшими} вероятностями.
3. Создать для них родительский узел, вес которого равен сумме весов этих двух узлов.
4. Заменить два выбранных узла на новый родительский узел в списке свободных узлов.
5. Повторять шаги 2–4, пока не останется один узел — \textbf{корень дерева}.

\subsection*{Дерево Хаффмана}

Результатом работы алгоритма является бинарное дерево, называемое \textbf{деревом Хаффмана}.
\begin{itemize}
    \item \textbf{Листья дерева} — это исходные символы алфавита.
    \item \textbf{Ветви} помечаются символами «0» (обычно левое ребро) и «1» (правое).
    \item \textbf{Код символа} — это последовательность меток на пути от корня к соответствующему листу.
\end{itemize}

\subsection*{Сравнение и пример}

\paragraph{Пример:} Алфавит $\{A: 0.4, B: 0.3, C: 0.2, D: 0.1\}$.
\begin{itemize}
    \item \textbf{Хаффман:} Складываем $D(0.1)$ и $C(0.2) \to CD(0.3)$. Затем $CD(0.3)$ и $B(0.3) \to BCD(0.6)$. Наконец $BCD(0.6)$ и $A(0.4) \to Root(1.0)$.
    \item \textbf{Коды:} $A$ получит короткий код «0» (1 бит), а $D$ — длинный «111» (3 бита).
\end{itemize}

\paragraph{Преимущества Хаффмана:}
1. \textbf{Минимальная избыточность:} средняя длина сообщения ближе всего к энтропии источника.
2. \textbf{Префиксность:} сообщение можно декодировать «на лету» без разделителей.
3. \textbf{Эффективность:} Хаффман всегда лучше или равен методу Фано.